\documentclass[aspectratio=169,11pt]{beamer}

\usetheme{Madrid}
\usecolortheme{default}
\usefonttheme{professionalfonts}
\setbeamertemplate{navigation symbols}{}
\setbeamertemplate{footline}[frame number]

\usepackage{amsmath, amssymb, booktabs}

\title{Crime, Safety, Environment}
\subtitle{And The Urban Feasible Set For Work}
\author{Special Topic 4: Urban and Labor -- Week 4}
\date{}

\begin{document}

\begin{frame}
  \titlepage
\end{frame}

\begin{frame}{Central question}
\begin{itemize}
    \item How do safety and environmental conditions change which jobs workers can take and how much work is worth?
    \item This is a labor lecture because risk changes search, commuting, hours, productivity, retention, sorting, and welfare.
    \item The relationship is two-sided: urban risk shapes work, and labor-market opportunity shapes crime and risky activity.
    \item Wages are not enough when risk and exposure differ across workers and places.
\end{itemize}
\end{frame}

\begin{frame}{Where Week 4 fits}
\begin{itemize}
    \item Week 1: workers choose job-residence bundles, not placeless jobs.
    \item Week 2: housing and moving costs shape real access.
    \item Week 3: segregation and neighborhoods create unequal job access.
    \item Week 4: crime, safety, pollution, heat, and noise alter the feasible set for work.
\end{itemize}
\end{frame}

\begin{frame}{Risk-adjusted value of work}
\[
V_{ijr\ell t}
=
w_{j\ell t}
- R_{rt}
- \tau_{ijr\ell t}
- \rho^S_{ijr\ell t}
- \rho^E_{ijr\ell t}
+ A_{rt}
+ \varepsilon_{ijr\ell t}
\]
\small
\begin{tabular}{p{0.20\linewidth} p{0.64\linewidth}}
\toprule
Term & Interpretation \\
\midrule
$w$ & wage at the job or workplace \\
$R$ & housing cost tied to residence \\
$\tau$ & generalized commute and schedule cost \\
$\rho^S$ & safety, harassment, violence, or public-space risk \\
$\rho^E$ & pollution, heat, noise, or other environmental exposure \\
$A$ & amenity value \\
\bottomrule
\end{tabular}
\end{frame}

\begin{frame}{Two-sided equilibrium framework}
\[
\Pr(C_{i\ell t}=1)
=
F\left(B_{i\ell t}-p_{\ell t}F_i-O_{i\ell t}+\eta_{i\ell t}\right)
\]
\small
\begin{tabular}{p{0.25\linewidth} p{0.58\linewidth}}
\toprule
Direction & Labor interpretation \\
\midrule
Risk $\rightarrow$ labor & safety and exposure change search, commutes, productivity, hours, retention \\
Labor $\rightarrow$ crime & job loss and weak earnings prospects lower the value of legal work \\
Risk $\rightarrow$ sorting & exposed workers and firms change residence, workplace, and sector choices \\
Risk $\rightarrow$ wages/rents & compensation and capitalization may be incomplete or unequal \\
Joint equilibrium & observed effects already include worker, firm, and housing responses \\
\bottomrule
\end{tabular}
\end{frame}

\begin{frame}{Risk and the feasible set}
\small
\begin{tabular}{p{0.26\linewidth} p{0.30\linewidth} p{0.28\linewidth}}
\toprule
Margin & Labor object & Example outcome \\
\midrule
Public safety & search radius, commute timing & employment, accepted wage, commute distance \\
Workplace harassment & retention, sector choice & turnover, segregation, inequality \\
Pollution, heat, noise & productivity, hours, attendance & output, fatigue, absenteeism \\
Environmental amenities & job and residence sorting & compensating differentials, welfare \\
Hidden harms & welfare versus wages & incomplete compensation \\
\bottomrule
\end{tabular}
\end{frame}

\begin{frame}{Crime, safety, and work}
\begin{itemize}
    \item Local labor-market opportunity changes the opportunity cost of crime and risky activity.
    \item Job displacement can move earnings, credit stress, and crime risk at the same time.
    \item Public-space risk changes commute routes, late shifts, job search radius, and schedule feasibility.
    \item Workplace harassment and violence affect quits, retention, occupational sorting, and inequality.
    \item The counterfactual must state whether safety, opportunity, or reporting changes.
\end{itemize}
\end{frame}

\begin{frame}{Environment and labor margins}
\small
\begin{tabular}{p{0.24\linewidth} p{0.35\linewidth} p{0.25\linewidth}}
\toprule
Exposure & Mechanism & Labor margin \\
\midrule
Pollution & health, fatigue, cognition & labor supply, absenteeism, output \\
Heat & physical stress and time allocation & hours, pace, task choice \\
Noise & attention and cognitive load & productivity, errors, fatigue \\
Amenities & sorting and compensation & wages, rents, real welfare \\
\bottomrule
\end{tabular}
\vspace{0.5em}
\begin{itemize}
    \item Exposure is unequal across residence, commute, workplace, task, and bargaining power.
    \item Adaptation can attenuate estimates while still leaving welfare losses.
\end{itemize}
\end{frame}

\begin{frame}{Why wages can mislead}
\begin{itemize}
    \item Risk may be hidden, poorly observed, or cumulative.
    \item Workers may lack mobility or bargaining power to demand compensation.
    \item Safer places can capitalize into rents rather than worker welfare.
    \item Firms can change staffing, tasks, or technology instead of wages.
    \item Equal wages can mask unequal commute risk, harassment, pollution, heat, and stress.
\end{itemize}
\end{frame}

\begin{frame}{Empirical designs and data}
\small
\begin{tabular}{p{0.25\linewidth} p{0.34\linewidth} p{0.26\linewidth}}
\toprule
Topic & Typical data & Design family \\
\midrule
Crime and opportunity & arrests, layoffs, matched employer data & layoff and demand shocks \\
Pollution and supply & monitors, satellites, worker panels & plant closures, inversions, panels \\
Workplace safety & surveys, HR records, complaints & event studies, within-firm designs \\
Heat, noise, output & weather, noise, task output & high-frequency panel designs \\
Commute safety & route, transit, crime data & policy and route shocks \\
\bottomrule
\end{tabular}
\end{frame}

\begin{frame}{Research spine}
\begin{itemize}
    \item Labor opportunity and crime: local labor markets, disadvantaged youth, job loss, and credit stress.
    \item Safety as a work constraint: crime timing, street harassment, workplace harassment, sexual violence.
    \item Environment and productivity: pollution, heat, noise, labor supply, and output.
    \item Welfare interpretation: compensation is often incomplete, and sorting can hide harms.
\end{itemize}
\end{frame}

\begin{frame}{Week 4 lab}
\begin{itemize}
    \item Reproduce: synthetic displacement event study for earnings, opportunity, credit stress, and crime risk.
    \item Diagnose: classify mechanisms as risk-to-labor, labor-to-crime, sorting, welfare, or equilibrium.
    \item Transfer: apply the framework to heat, pollution, noise, and worker productivity.
    \item Extension: add a commute-safety shock that changes feasible shifts while holding wages fixed.
\end{itemize}
\end{frame}

\begin{frame}{Takeaways}
\begin{itemize}
    \item Crime, safety, and environment are labor-market constraints when they alter feasible work.
    \item Labor-market opportunity also shapes crime and risky activity.
    \item Environment matters through labor supply, productivity, health, cognition, and adaptation.
    \item The best papers name the margin, identify the counterfactual, and avoid treating wages as welfare.
\end{itemize}
\end{frame}

\end{document}
